% !TEX root = ../journal_0421052099.tex
% ---------------------------------------------------------------------
% No new numbers, no new claims.
%
% Close on the attribution argument, not on the system. Since the
% September 2026 reframe the paper IS the attribution argument, so the
% conclusion restates what was measured, what it was attributed to, and
% what the measurement could not see. The thesis's closing remark --
% fact verification is the easier half, relevance verification is the
% half still open -- sits in the second paragraph.
%
% Trimmed 2026-09-23, from ~500 words to ~250. This section carries no
% number and no \Cref, so it is the one place in the manuscript where
% cutting costs nothing but restatement: four paragraphs became two,
% and each finding is now named once rather than re-argued.
% ---------------------------------------------------------------------

\section{Conclusion}\label{sec:conclusion}

We built an agentic knowledge-graph question-answering system whose four
components can be removed one at a time, ran it against four baselines
on one frozen backbone under one call budget, and took it apart. The
aggregate result is the one the field usually reports: the agentic
system has the highest accuracy in every column, at roughly half the
model calls of the agentic baseline. The rest of the paper is what that
result turns out to be made of. The entire margin over the agentic
baseline comes from the questions on which the shared budget truncates
it, so it is a finding about affordability under a fixed spend rather
than about reasoning power. Of four components, three could not be shown
to earn their cost at the power available --- including the
claim-verification layer the system was designed around, which changes
the answer on one question per dataset and delivers an output contract
rather than an accuracy gain. The fourth was actively harmful on the
shallower benchmark: removing the planner improved accuracy and cut
token use, because decomposing a question that has no compositional
structure discards the context that would have answered it.

The shape of the residue matters as much as the margin. A system that
traverses a real graph and copies its answers out of real triples can be
made to assert nothing ungrounded, and the navigation paradigm achieves
that, not any checking layer built on top. What remains is an answer
that is grounded, connected, and true, and is still not the answer to
the question that was asked; systems sharing a graph and a backbone
converge on that answer together, so a protocol that takes their
agreement as evidence ratifies it. Fact verification against a knowledge
graph turns out to be the easier half of the problem, and relevance
verification is the half still open. The attribution result is finally a
claim about how this field measures rather than about this architecture.
A component adopted because the assembly containing it performed well
may be contributing nothing, or less than nothing, and a margin reported
as an architectural advantage may be a budget regime. The only way to
know is to remove the component, vary the budget, and test.
